% Part: proof-theory
% Chapter: sequent-calculus
% Section: admissible-derivable

\documentclass[../../../include/open-logic-section]{subfiles}

\begin{document}

\olfileid{pt}{seq}{adm}

\olsection{قواعد پذیرفتنی و مشتق‌شدنی}

بسیاری از نتایج نظریهٔ برهان به این پرسش مربوط‌اند ــ یا از نتایجی دربارهٔ
آن استفاده می‌کنند ــ که آیا هر آنچه با قاعده‌ای معین می‌توان !!{prove} کرد،
بدون آن قاعده نیز !!{prove}شدنی است یا نه. قضیهٔ حذف برش شناخته‌شده‌ترین
نمونه است. این قضیه نشان می‌دهد هر دنباله‌ای را که با قواعد یک دستگاه
دنباله‌ایِ شامل~\Cut{} بتوان !!{prove} کرد، بدون~\Cut نیز می‌توان
!!{prove} کرد. چون این‌گونه ویژگی بسیار مهم است، نامی ویژه دارد.

\begin{defn}
قاعدهٔ~$R$ در یک دستگاه دنباله‌ای \emph{پذیرفتنی} است هرگاه هر
دنباله‌ای را که در آن دستگاه همراه با قاعدهٔ~$R$ بتوان اثبات کرد، بدون~$R$
نیز بتوان اثبات کرد.
\end{defn}

\begin{ex}\ollabel{ex:land-deriv}
قاعدهٔ~$\LeftR{\land}$ در~\Log{G1c} و~\Log{G3c} صورت‌های متفاوتی دارد.
در~\Log{G1c} دو نسخه از قاعده هست: یکی با مؤلفهٔ عطفی چپ~$!A$ از
!!{formula} اصلی~$!A \land !B$ در مقدمه، و دیگری با مؤلفهٔ عطفی راست~$!B$:
\[
\Axiom$ !A, \Gamma \fCenter \Delta$
\RightLabel{$\LeftR{\land}_1$}
\UnaryInf$ !A \land !B, \Gamma \fCenter \Delta$
\DisplayProof
\qquad
\Axiom$!B, \Gamma \fCenter \Delta$
\RightLabel{$\LeftR{\land}_2$}
\UnaryInf$!A \land !B, \Gamma \fCenter \Delta$
\DisplayProof
\]
در مقابل،~\Log{G3c} تنها یک قاعده دارد:
\[
\Axiom$ !A, !B, \Gamma \fCenter \Delta$
\RightLabel{$\LeftR{\land}_3$}
\UnaryInf$ !A \land !B, \Gamma \fCenter \Delta$
\DisplayProof
\]
اکنون می‌توان پرسید: «آیا قاعدهٔ~$\LeftR{\land}_3$ در~\Log{G3c}،
در~\Log{G1c} پذیرفتنی است؟» پاسخ مثبت است: هر استنتاج با
$\LeftR{\land}_3$ را می‌توان مستقیماً و تنها با قواعد~\Log{G1c} شبیه‌سازی
کرد. افزون بر~$\LeftR{\land}_1$ و~$\LeftR{\land}_2$، به قاعدهٔ ساختاری
$\LeftR{\Contraction}$ نیز نیاز داریم:
\[
\Axiom$!A, !B, \Gamma \fCenter \Delta$
\RightLabel{$\LeftR{\land}_1$}
\UnaryInf$!A \land !B, !B, \Gamma \fCenter \Delta$
\RightLabel{$\LeftR{\land}_2$}
\UnaryInf$!A \land !B, !A \land !B, \Gamma \fCenter \Delta$
\RightLabel{\LeftR{\Contraction}}
\UnaryInf$!A \land !B, \Gamma \fCenter \Delta$
\DisplayProof
\]
\end{ex}

شبیه‌سازیِ استنتاجی با یک قاعده به‌وسیلهٔ قواعد دیگر، یکی از راه‌های نشان‌دادن
پذیرفتنی‌بودن قاعده است؛ اما تنها راه نیست و برخی قواعد پذیرفتنی‌اند، بی‌آنکه
چنین شبیه‌سازی‌ای برایشان ممکن باشد. قواعدی را که می‌توان بدین‌سان شبیه‌سازی
کرد \emph{مشتق‌شدنی} می‌نامند.

\begin{defn}
قاعدهٔ~$R$،
\[
\AxiomC{$S_1$}
\AxiomC{$\dots$}
\AxiomC{$S_n$}
\RightLabel{$R$}
\TrinaryInfC{$S$}
\DisplayProof
\]
در یک دستگاه \emph{مشتق‌شدنی} نامیده می‌شود هرگاه !!{proof}
طرح‌واره‌ای از دنبالهٔ~$S$ وجود داشته باشد که در آن، قواعد و اصول موضوعهٔ
دستگاه همراه با دنباله‌های آغازیِ اضافی به صورت‌های~$S_1$، \dots،~$S_n$
به کار رفته باشند.
\end{defn}

!!{proof} داده‌شده در~\olref{ex:land-deriv} نشان می‌دهد
$\LeftR{\land}_3$ قاعده‌ای مشتق‌شدنی در~\Log{G1c} است، زیرا آن
!!{proof}، !!{proof} طرح‌واره‌ای از نتیجهٔ~$\LeftR{\land}_3$ بر پایهٔ
مقدمات~$\LeftR{\land}_3$ به‌عنوان دنباله‌های آغازی و با استفادهٔ صرف از
قواعد~\Log{G1c} است. (هیچ‌یک از اصول موضوعهٔ~\Log{G1c} در آن به کار
نرفته است.)

\begin{prop}\ollabel{prop:lor-G3-adm}
قواعد~$\LeftR{\land}$ و~$\RightR{\lor}$ در~\Log{G3c}، در~\Log{G1c}
مشتق‌شدنی‌اند.
\end{prop}

\begin{proof}
  برهان~$\LeftR{\land}$ در~\olref{ex:land-deriv} داده شد. برهان
  $\RightR{\lor}$ به‌عنوان تمرین واگذار می‌شود.
\end{proof}

\begin{prob}
نشان دهید قاعدهٔ~$\RightR{\lor}$ در~\Log{G3c}، در~\Log{G1c}
مشتق‌شدنی است.
\end{prob}

\begin{prob}
فرض کنید~$\liff$ !!{operator} تعریف‌شده باشد: $!A \liff !B$ کوتاه‌شدهٔ
$(!A \lif !B) \land (!B \lif !A)$ است. نشان دهید قواعد
\begin{enumerate}
\item \Axiom$\Gamma, !A, !B \fCenter \Delta$
\Axiom$\Gamma \fCenter \Delta, !A, !B$
\RightLabel{\LeftR{\liff}}
\BinaryInf$!A \liff !B, \Gamma \fCenter \Delta$
\DisplayProof
\item
\Axiom$!A, \Gamma \fCenter \Delta, !B$
\Axiom$!B, \Gamma \fCenter \Delta, !A$
\RightLabel{\RightR{\liff}}
\BinaryInf$\Gamma \fCenter \Delta, !A \liff !B$
\DisplayProof
\end{enumerate}
در (الف)~\Log{G1c} و (ب)~\Log{G3c} مشتق‌شدنی‌اند.
\end{prob}


\begin{ex}\ollabel{ex:weak-adm}
قاعدهٔ~$\RightR{\Weakening}$ در~\Log{G1c}،
\[
\Axiom$ \Gamma \fCenter \Delta$
\RightLabel{\RightR{\Weakening}}
\UnaryInf$ \Gamma \fCenter \Delta, !A$
\DisplayProof
\]
در~\Log{G3c} پذیرفتنی است. اندیشهٔ برهان ساده است: فرض کنید در~\Log{G3c}
!!a{proof} به نام~$\pi$ از مقدمهٔ~$\Gamma \Sequent \Delta$ وجود دارد.
!!{formula}~$!A$ را به تالیِ هر دنباله در~$\pi$ بیفزایید. اصول همچنان اصل
می‌مانند و استنتاج‌های درست همچنان درست‌اند؛ همهٔ !!{formula}های پیشین نیز
حفظ می‌شوند. دنبالهٔ پایانیِ !!{proof} جدید
$\Gamma \Sequent \Delta, !A$ است.

بااین‌حال، قاعدهٔ~$\RightR{\Weakening}$ در~\Log{G3c} مشتق‌شدنی نیست.
زیرا فرض کنید چنین بود؛ یعنی می‌شد با اصول و قواعد~\Log{G3c}، و احتمالاً
دنبالهٔ آغازیِ اضافی~$\Gamma \Sequent \Delta$، !!a{proof} از
$\Gamma \Sequent \Delta, !A$ یافت. اگر این !!{proof}، آن‌گونه که برای
مشتق‌شدنی‌بودن~$\RightR{\Weakening}$ لازم است، کاملاً طرح‌واره‌ای باشد،
با حذف~$\Gamma$ و~$\Delta$ می‌توان آن را به !!a{proof} از~$\Sequent !A$
تبدیل کرد. در این صورت دنباله‌های آغازیِ اضافی~$\Gamma \Sequent \Delta$
به دنبالهٔ تهی~$\quad \Sequent \quad$ بدل می‌شوند. چون دنبالهٔ تهی اصلاً
هیچ !!{formula}ی ندارد، ولی همهٔ قواعد~\Log{G3c} دست‌کم یک !!{formula}
فعال در مقدماتشان می‌طلبند، دنبالهٔ تهی نمی‌تواند مقدمهٔ هیچ استنتاجی در این
!!{proof} طرح‌واره‌ای باشد. در نتیجه، این !!{proof} تنها در صورتی
طرح‌واره‌ای است که در واقع دنبالهٔ اضافی~$\Gamma \Sequent \Delta$ را
به‌عنوان دنبالهٔ آغازی به کار نبرد. به بیان دیگر، اشتقاقی طرح‌واره‌ای از
$\Gamma \Sequent \Delta, !A$ تنها با اصول و قواعد~\Log{G3c} می‌بود و
نشان می‌دادیم هر دنباله‌ای به این صورت در~\Log{G3c} !!a{proof} دارد. اما
این ناممکن است، زیرا~\Log{G3c} درست است و فقط دنباله‌هایی را می‌تواند
!!{prove} کند که در هر !!{structure} صادق‌اند. با گرفتن
$\Gamma = \Delta = \emptyset$ و~$!A \ident \lfalse$، دستگاه~\Log{G3c}
ناگزیر می‌بود، برای مثال، دنبالهٔ~$\quad \Sequent \lfalse$ را
!!{prove} کند، حال آنکه چنین نمی‌کند.
\end{ex}

اکنون برهانی استقرایی و خوش‌ساخت برای پذیرفتنی‌بودن تضعیف در~\Log{G3c}
می‌دهیم. در واقع، همان‌گونه که استدلال شهودی نشان می‌دهد، نتیجه‌ای اندکی
قوی‌تر برقرار است: افزودن~$!A$ به سمت راست هر دنبالهٔ~$\pi$ ساختار
!!{proof} و، به‌ویژه، ارتفاع آن را تغییر نمی‌دهد. چون این نکته مهم است،
تعریفی ویژه می‌طلبد.

\begin{defn}
قاعدهٔ~$R$،
\[
\AxiomC{$S_1$}
\AxiomC{$\dots$}
\AxiomC{$S_n$}
\RightLabel{$R$}
\TrinaryInfC{$S$}
\DisplayProof
\]
در یک دستگاه \emph{پذیرفتنی با حفظ !!{height}} (یا
\emph{!!{hp}-پذیرفتنی}) است هرگاه، اگر~$\Proves[h_i] S_i$ برای
$i = 1$، \dots،~$n$، آنگاه~$\Proves[m] S$ برای عددی~$m$ چنان‌که
$m \le \max(h_1, \dots, h_n)$.
\end{defn}

برای پذیرفتنی با حفظ !!{height} بودنِ قاعده کافی است هرگاه مقدمات~$S_i$
دارای !!{proof}s~$\pi_i$ با !!{height}~$h_i = \pheight{\pi_i}$ باشند،
!!a{proof}~$\pi$ از نتیجه با !!{height}~$\pheight{\pi}$ وجود داشته باشد
که این ارتفاع از بیشینهٔ
$h_i$ها بیشتر نباشد. به‌ویژه، اگر تنها یک مقدمهٔ~$S_1$ باشد، کافی است
$\pheight{\pi} \le \pheight{\pi_1}$. در مورد
$\RightR{\Weakening}$ و~$\LeftR{\Weakening}$ حتی
$\pheight{\pi} = \pheight{\pi_1}$ داریم، اما برابری لازم نیست.

\begin{prop}\ollabel{prop:weak-G3c-adm}
قواعد~$\RightR{\Weakening}$ و~$\LeftR{\Weakening}$ در~\Log{G1c}،
در~\Log{G3c} با حفظ !!{height} پذیرفتنی‌اند.
\end{prop}

\begin{proof}
  باید نشان دهیم اگر~$\Log{G3c} \Proves \Gamma \Sequent \Delta$ با
  !!{proof}~$\pi$ از !!{height}~$\pheight{\pi} = n$، آنگاه در
  ~\Log{G3c} !!{proof}~$\pi'$ از~$\Gamma \Sequent \Delta, !A$ وجود
  دارد که~$\pheight{\pi'} = n$. بر~$n$ استقرا می‌کنیم. اگر~$n=0$، آنگاه
  $\pi$ یک اصل~$\Gamma \Sequent \Delta$ است: یا یک~$!C$ اتمی هم در
  $\Gamma$ و هم در~$\Delta$ رخ می‌دهد، یا دنباله اصلی از نوع اصلِ
  $\lfalse$ در چپ است. در هر دو حالت~$\Gamma \Sequent \Delta, !A$ نیز
  اصل است. اگر~$n = \pheight{\pi} >0$، !!{proof}~$\pi$ استنتاج پایانی
  دارد. بر حسب قاعدهٔ به‌کاررفته حالت‌بندی می‌کنیم. تنها یک نمونه را
  می‌آوریم: فرض کنید استنتاج پایانی~$\RightR{\land}$ باشد. آنگاه
  $\Delta = \Delta', !B \land !C$ و زیر-!!{proof}s بی‌واسطهٔ~$\pi_1$
  و~$\pi_2$ در آن استنتاج، به‌ترتیب دنباله‌های پایانی
  $\Gamma \Sequent \Delta', !B$ و~$\Gamma \Sequent \Delta', !C$ را
  دارند. یعنی~$\pi$ چنین است:
  \[
  \AxiomC{}
  \RightLabel{$\pi_1$}
  \Deduce$\Gamma \fCenter \Delta', !B$
  \AxiomC{}
  \RightLabel{$\pi_2$}
  \Deduce$\Gamma \fCenter \Delta', !C$
  \RightLabel{\RightR{\land}}
  \BinaryInf$\Gamma \fCenter \Delta', !B \land !C$
  \DisplayProof
  \]
  همچنین~$n = \max(\pheight{\pi_1}, \pheight{\pi_2}) + 1$. پس
  $\pheight{\pi_1} < n$ و~$\pheight{\pi_2} < n$ و فرض استقرا اعمال
  می‌شود: !!{proof}s~$\pi_1'$ و~$\pi_2'$ به‌ترتیب از
  $\Gamma \Sequent \Delta', !B, !A$ و
  $\Gamma \Sequent \Delta', !C, !A$ وجود دارند. با قاعدهٔ
  $\RightR{\land}$، !!a{proof}~$\pi'$ از
  $\Gamma \Sequent \Delta', !B \land !C, !A$ به دست می‌آوریم:
  \[
  \AxiomC{}
  \RightLabel{$\pi_1'$}
  \Deduce$\Gamma \fCenter \Delta', !B, !A$
  \AxiomC{}
  \RightLabel{$\pi_2'$}
  \Deduce$\Gamma \fCenter \Delta', !C, !A$
  \RightLabel{\RightR{\land}}
  \BinaryInf$\Gamma \fCenter \Delta', !B \land !C, !A$
  \DisplayProof
  \]
  داریم
  $\pheight{\pi'} =
  \max(\pheight{\pi_1'}, \pheight{\pi_2'}) + 1 \le
  \max(\pheight{\pi_1}, \pheight{\pi_2}) + 1 = \pheight{\pi}$.
  حالت‌های دیگرِ قاعدهٔ پایانی به همین روش‌اند، و برهان تضعیف چپ نیز
  متقارن است.
\end{proof}

\olref{prop:weak-G3c-adm} بسیار سودمند است. برای کاربرد مکرر آن نمادی
معرفی می‌کنیم: اگر~$\pi$ !!a{proof} از~$\Gamma \Sequent \Delta$ باشد،
$\pi[\Pi \Sequent \Lambda]$ حاصلِ اعمال پذیرفتنی‌بودن تضعیف چپ با همهٔ
!!{formula}s موجود در~$\Pi$ و پذیرفتنی‌بودن تضعیف راست با همهٔ
!!{formula}s موجود در~$\Lambda$ است. این !!a{proof} از
$\Gamma, \Pi \Sequent \Delta, \Lambda$ است.

\end{document}
